% Copy this file when starting a new chapter.
% Replace the placeholders with real content; do not \include this file directly.
% If a chapter needs an explicit exception under CONTENT_README.md, record it
% immediately below \chapter{...} as a comment.
% DELETE ALL "% Layout note:" and "% Pattern note:" blocks before committing a
% real chapter file copied from this template.

\chapter{Chapter Title}

This chapter develops the main ideas of the chapter in one or two opening
paragraphs. The overview should identify the mathematical territory, connect
the chapter to earlier material, and preview the central results students
should look forward to.

Use the opening overview for motivation and orientation only. Do not introduce
new notation here unless it is unavoidable.

\section{Section Title}

Open the section with a short paragraph that explains what problem this
section addresses and why it matters.

% Layout note:
% - section and subsection headings now reserve space automatically, so avoid
%   manual \newpage/\pagebreak just to keep a heading with its first lines
% - long titles may wrap naturally; the shared heading style will align them

\subsection{Subsection title}

Continue with textbook-style prose before the first formal environment.

\begin{definition}
State the new concept cleanly and without examples or side remarks.
\end{definition}

\begin{remark}
Use a remark for interpretation, contrast, or a brief warning if it helps.
\end{remark}

\begin{workedexample}
\begin{example}
Evaluate \(\lim_{x\to a^-}f(x)\), \(\lim_{x\to a^+}f(x)\), and determine
whether \(\lim_{x\to a}f(x)\) exists.
\end{example}

\begin{solution}
If \(x\to a^{-}\), then the relevant branch or sign information gives
\(\lim_{x\to a^-}f(x)=L_-\).

If \(x\to a^{+}\), then the relevant branch or sign information gives
\(\lim_{x\to a^+}f(x)=L_+\).

Since \(L_- \ne L_+\), \(\lim_{x\to a}f(x)\) does not exist.
\end{solution}
\end{workedexample}

% Pattern note:
% - keep `Solution.` inline unless the body opens with a block; then start with
%   `\solutionbreak`
% - if a solution ends with display math, place `\qedhere` on the last display
%   line so the closing box stays attached to the mathematics
% - use one local display grammar per algebraic unit instead of mixing centered
%   displays, aligned blocks, and prose-side notes
% - `aligneddisplay`: short top-to-bottom peers or derivation lines
% - `conditiondisplay`: formula plus a domain/range/branch condition
% - `\pairdisplay{...}{...}`: exactly two short formulas meant to be compared
% - `\iffwithconditions{...}{...}{...}`: equivalence plus one grouped
%   restriction set (brace is semantic)
% - `\iffstackeddisplay{...}{...}{...}`: equivalence plus two equal-status
%   stacked criteria
% - if an example has no worked solution, use `example` alone rather than
%   `workedexample`
% - see CONTENT_README.md for the full policy before committing

\begin{aligneddisplay}
u(x) &= x^2-1,\\
v(x) &= \sqrt{x+3}.
\end{aligneddisplay}

\begin{conditiondisplay}
f(x) &= x, \qquad & 0 \le x \le 1,\\
g(x) &= x^2, \qquad & -1 \le x \le 1.
\end{conditiondisplay}

\pairdisplay{\lim_{x\to a^-}f(x)=L,}{\lim_{x\to a^+}f(x)=L.}

\iffwithconditions{f^{-1}(y)=x}{f(x)=y,}{y\in B.}

\iffstackeddisplay{\lim_{x\to a}f(x)=L}{\lim_{x\to a^-}f(x)=L,}{\lim_{x\to a^+}f(x)=L.}

\begin{aligneddisplay}
f(x+y) &= f(x)+f(y),\\
f(cx) &= cf(x).
\end{aligneddisplay}

\begin{workedexample}
\begin{example}
Evaluate
\[
\lim_{x\to a}\frac{\sqrt{x+b}-c}{x-a}.
\]
\end{example}

\begin{solution}
Direct substitution gives the indeterminate form \(\frac{0}{0}\), so we
rationalize:
\[
\frac{\sqrt{x+b}-c}{x-a}
\cdot
\frac{\sqrt{x+b}+c}{\sqrt{x+b}+c}
=
\frac{x+b-c^2}{(x-a)\bigl(\sqrt{x+b}+c\bigr)}.
\]
After simplification, we obtain
\[
\frac{1}{\sqrt{x+b}+c}.
\]
Therefore,
\[
\lim_{x\to a}\frac{\sqrt{x+b}-c}{x-a}
=
\frac{1}{2c}.
\]
\end{solution}
\end{workedexample}

\begin{example}
State a brief illustrative example here when no worked solution follows.
\end{example}

\begin{theorem}
State a formal result here when the section needs one.
\end{theorem}

\begin{proof}
Give a genuine proof only when the book's level and the manuscript warrant it.
\end{proof}

% TODO: add \subsection*{Exercises} block with end-of-section problems for
% this section, following CONTENT_README.md.
